\documentclass[11pt]{article}

\usepackage[margin=1in]{geometry}
\usepackage{amsmath}
\usepackage{amssymb}
\usepackage[hidelinks]{hyperref}

\title{Better flows for the hat layout:\\incidence, rigidity and
Laplacian solves}
\author{}
\date{5 October 2026}

\begin{document}
\maketitle

Yes. I think there is a substantially faster formulation here, and the
object you want to exploit is actually \textbf{the incidence/rigidity
matrix even more than the adjacency matrix}.

Your current optimization sounds like it is doing explicit time stepping on
a stiff geometric energy. A deformation at one end of the fabric only
propagates a little way per iteration, which is exactly the situation where
gradient descent takes thousands of iterations. Instead, you can linearize
the local nonlinear geometry and solve the resulting \textbf{graph-wide
sparse system} at every outer iteration. That lets information propagate
across the entire knit in one solve.

The resulting method sits very naturally between \textbf{stress
majorization / projective dynamics}, \textbf{rigidity theory}, and a kind
of \textbf{prescribed-curvature discrete Ricci flow}. Stress majorization
is already known to turn nonlinear graph-distance objectives into repeated
Laplacian solves, and projective dynamics uses essentially the same
local/global idea for geometric constraints.
\href{https://www.researchgate.net/publication/30508817_Graph_Drawing_by_Stress_Majorization?utm_source=chatgpt.com}{[ResearchGate]}

\section*{1. Put the knitting graph into incidence-matrix form}

Let
\[
G=(V,E),\qquad |V|=n,\quad |E|=m,
\]
and let
\[
X=
\begin{bmatrix}
x_1^\top\\
\vdots\\
x_n^\top
\end{bmatrix}\in\mathbb R^{n\times 3}
\]
be your vertex positions.

Choose an arbitrary orientation for the edges and construct the signed
incidence matrix
\[
B\in\mathbb R^{m\times n}.
\]
Then
\[
Z=BX\in\mathbb R^{m\times3}
\]
contains all edge vectors. For edge \(e=(i,j)\),
\[
z_e=x_i-x_j,\qquad
r_e=\|z_e\|,\qquad
u_e=\frac{z_e}{r_e}.
\]
Your weighted Laplacian is
\[
L_W=B^\top W B.
\]
This representation is almost tailor-made for your problem.

Suppose, approximately, your current energy is
\[
E(X)=
\lambda_\ell E_\ell(X)
+\lambda_\Omega E_\Omega(X)
+\lambda_R E_R(X),
\]
where \(E_\ell\) is edge-length distortion, \(E_\Omega\) is your
solid-angle term, and \(E_R\) is universal repulsion.

\bigskip\hrule

\section*{2. Your log edge-length loss already has a hidden graph Laplacian}

Suppose your edge term is
\[
E_\ell(X)
=
\frac12\sum_e
w_e
\left[
\log\frac{r_e}{\ell_e}
\right]^2,
\]
where \(\ell_e\) is the desired yarn/stitch edge length.

Define
\[
s_e=\log(r_e/\ell_e).
\]
Then
\[
\frac{\partial E_\ell}{\partial z_e}
=
w_e\,s_e\,\frac{z_e}{r_e^2}.
\]
Consequently,
\[
\boxed{
\nabla_X E_\ell
=
B^\top
\operatorname{diag}
\left(
\frac{w_e s_e}{r_e^2}
\right)
BX.
}
\]
So even the exact gradient is already a \textbf{weighted graph-Laplacian
operation}.

The problem is that doing
\[
X^{k+1}=X^k-\eta\nabla E
\]
uses that Laplacian explicitly rather than inversely.

That distinction is enormous.

If
\[
Lq_i=\lambda_iq_i,
\]
ordinary gradient descent changes a mode by roughly
\[
c_i^{k+1}
=
(1-\eta\lambda_i)c_i^k.
\]
Because stability forces \(\eta\) to accommodate \(\lambda_{\max}\), the
modes near \(\lambda_2\) move extremely slowly. Roughly speaking, the
convergence rate is controlled by
\[
\kappa(L)=\frac{\lambda_{\max}}{\lambda_2}.
\]
That looks very much like the ``gradually changes shape over 10k
iterations'' behavior you're seeing.

Instead, you want to approximately apply
\[
L^\dagger
\]
to the force.

Then a disturbance at one vertex can affect the whole connected component
\textbf{in one iteration}.

\bigskip\hrule

\section*{3. Even better: the proper matrix is the rigidity matrix}

For an edge \(e=(i,j)\), define the usual rigidity-matrix row
\[
R_e(X)
=
[
0,\ldots,
z_e^\top,
\ldots,
-z_e^\top,
\ldots,0
].
\]
The rigidity matrix is the Jacobian of squared edge lengths; this is
precisely the standard object from framework rigidity theory.
\href{https://link.springer.com/article/10.1007/s00454-021-00275-7?utm_source=chatgpt.com}{[Springer]}

For your log-length residual,
\[
s_e=\log(r_e/\ell_e),
\]
we have
\[
\nabla s_e
=
\frac{1}{r_e^2}R_e(X).
\]
Therefore
\[
J_\ell
=
D_{r^{-2}} R(X),
\]
and its Gauss--Newton Hessian is
\[
\boxed{
H_\ell^{GN}
=
R^\top
D_{w_er_e^{-4}}
R.
}
\]
Equivalently, each edge contributes the \(3\times3\) stiffness
\[
\frac{w_e}{r_e^2}u_eu_e^\top
\]
between its two incident vertices.

This is a \textbf{vector-valued graph Laplacian / framework stiffness
matrix}.

That is probably the matrix you actually want to solve against.

Instead of
\[
X_{k+1}=X_k-\eta g_k,
\]
do
\[
\boxed{
(H_k+\mu P)\Delta X=-g_k
}
\]
and
\[
X_{k+1}=X_k+\alpha\Delta X.
\]
Here \(P\) can be something like
\[
P=L_0\otimes I_3
\]
rather than ordinary \(I\). This gives you a very natural
\emph{graph-Sobolev} Levenberg--Marquardt method.

\bigskip\hrule

\section*{4. There is an even cheaper local/global formulation}

For edge preservation, I would seriously test this before implementing a
full Newton method.

Consider ordinary edge stress
\[
E_s(X)
=
\frac12
\sum_e
w_e(\|z_e\|-\ell_e)^2.
\]
Introduce a vector \(p_e\) constrained to lie on the sphere
\[
\|p_e\|=\ell_e.
\]
Then
\[
\boxed{
(\|z_e\|-\ell_e)^2
=
\min_{\|p_e\|=\ell_e}
\|z_e-p_e\|^2.
}
\]
So the nonlinear problem becomes
\[
\min_{X,P}
\frac12
\sum_e
w_e\|(BX)_e-p_e\|^2,
\qquad
\|p_e\|=\ell_e.
\]
Now alternate.

The \textbf{local step} is trivial:
\[
p_e
=
\ell_e
\frac{z_e}{\|z_e\|}.
\]
The \textbf{global step} is
\[
\min_X
\frac12\|W^{1/2}(BX-P)\|_F^2,
\]
giving
\[
\boxed{
B^\top W B\,X
=
B^\top W P.
}
\]
Or simply
\[
\boxed{
L_W X=B^\top WP.
}
\]
That's three sparse Laplacian solves.

And, crucially, if \(W\) is constant, \textbf{you factor \(L_W\) once} and
reuse it forever.

This is basically the same local/global pattern underlying stress
majorization, ARAP methods, and projective dynamics. Stress majorization was
introduced to graph drawing specifically because it provides faster and
more stable optimization than small local steps; projective dynamics
similarly exploits independent local projections plus a constant global
system.
\href{https://www.researchgate.net/publication/30508817_Graph_Drawing_by_Stress_Majorization?utm_source=chatgpt.com}{[ResearchGate]}

For your knitting application, it is almost suspiciously well matched.

\bigskip\hrule

\section*{5. You don't have to give up the log-length loss}

You can turn the log loss into an iteratively reweighted version of the same
Laplacian solve.

We want a surrogate
\[
\frac12\alpha_e^{(k)}
(r_e-\ell_e)^2
\]
whose derivative agrees with your log loss at the current \(r_e\).

Your log-loss derivative is
\[
\frac{d}{dr}
\frac12 w_e\log^2(r/\ell_e)
=
\frac{w_e\log(r/\ell_e)}{r}.
\]
The quadratic stress derivative is
\[
\alpha_e(r-\ell_e).
\]
Thus choose
\[
\boxed{
\alpha_e
=
w_e
\frac{\log(r_e/\ell_e)}
{r_e(r_e-\ell_e)}.
}
\]
This is positive on both sides of the desired length, and
\[
\lim_{r\rightarrow\ell}
\alpha_e
=
\frac{w_e}{\ell_e^2}.
\]
Then perform
\[
p_e=\ell_eu_e,
\]
followed by
\[
\boxed{
L_{\alpha} X^{k+1}
=
B^\top D_{\alpha}P.
}
\]
This is a lagged-weight quasi-Newton iteration.

You can use a line search against your \emph{actual} log energy, so you
retain the original objective while using a much better optimizer.

Near convergence,
\[
\alpha_e\simeq\frac{w_e}{\ell_e^2},
\]
so the matrix becomes essentially constant and you can reuse a
factorization or a very good preconditioner.

\bigskip\hrule

\section*{6. Now add the solid-angle term}

Suppose
\[
c_v(X)
=
\Omega_v(X)-\Omega_v^\star
\]
and
\[
E_\Omega
=
\frac12
\sum_v q_v c_v^2.
\]
The important thing is not even the exact solid-angle formula.

The important fact is:
\[
\Omega_v
=
\Omega_v\left(x_v,\{x_j:j\in N(v)\}\right).
\]
So its Jacobian has nonzeros only in the \textbf{one-ring around \(v\)}.

Let
\[
J_\Omega
=
\frac{\partial(c_1,\ldots,c_n)}
{\partial\operatorname{vec}X}.
\]
Then linearizing gives
\[
c(X+\Delta)
\simeq
c(X)+J_\Omega\Delta.
\]
Thus
\[
H_\Omega^{GN}
=
J_\Omega^\top W_\Omega J_\Omega.
\]
And here's where your adjacency-matrix-powers intuition becomes
particularly interesting:
\[
\boxed{
\operatorname{supp}
(J_\Omega^\top J_\Omega)
\subseteq
\operatorname{supp}\big((I+A)^2\big).
}
\]
Why?

Because two vertices interact in the angle Hessian precisely when they
occur in a common vertex star. They are therefore either neighbors or share
a neighbor.

So the solid-angle Hessian naturally has an
\[
A+A^2
\]
kind of sparsity pattern.

I would \textbf{not} add \(A^2\) manually.

Let the geometry derive the coefficients, and exploit the fact that its
sparsity is exactly the graph's two-hop sparsity.

\bigskip\hrule

\section*{7. One outer iteration can therefore be one sparse solve}

Using the local/global edge surrogate, linearized solid angle, and explicit
repulsion, define
\[
P_e^{(k)}
=
\ell_e
\frac{x_i^k-x_j^k}
{\|x_i^k-x_j^k\|}.
\]
Let
\[
K_\ell=L_{\alpha}\otimes I_3.
\]
The edge surrogate gradient at \(X^k\) is
\[
g_\ell
=
K_\ell x^k
-
\operatorname{vec}(B^\top D_\alpha P).
\]
For the angular term,
\[
g_\Omega
=
J_\Omega^\top W_\Omega c.
\]
Then solve
\[
\boxed{
\left[
K_\ell
+
\lambda_\Omega
J_\Omega^\top W_\Omega J_\Omega
+
\mu P
\right]\Delta x
=
-
\left[
g_\ell+
\lambda_\Omega g_\Omega+
\lambda_R g_R
\right].
}
\]
Then
\[
x^{k+1}=x^k+\alpha_{\rm line}\Delta x.
\]
That is the formulation I'd start with.

Notice what happened:
\[
\text{thousands of scalar-ish gradient steps}
\]
became
\[
\text{a few nonlinear outer iterations}
+
\text{sparse graph solves}.
\]
For a knit graph, \(m=O(n)\), and the matrices are sparse.

With PCG + algebraic multigrid, each solve can get close to \(O(m)\)
behavior in practice.

\bigskip\hrule

\section*{8. Treat universal repulsion differently}

Repulsion is the one component that is fundamentally \textbf{not
graph-local}.

If
\[
E_R=\sum_{i<j}\psi(\|x_i-x_j\|),
\]
its Hessian is dense.

I would therefore not put its full Hessian into the solve.

Compute
\[
g_R=\nabla E_R
\]
using a Barnes--Hut octree or FMM-like approximation, and put it on the
right-hand side.

So your iteration is basically
\[
\underbrace{
(\text{edge stiffness}
+
\text{angular stiffness})
}_{\text{implicit}}
\Delta x
=
-
\underbrace{
(\text{edge}+\text{angle}+\text{repulsion forces})
}_{\text{forces}}.
\]
This is analogous to a semi-implicit physical simulation: treat the stiff
local interactions implicitly and the long-range force explicitly.

That alone can permit vastly larger effective step sizes.

\bigskip\hrule

\section*{9. This also suggests a very nice ADMM/projective formulation}

There is another formulation that may fit your implementation particularly
well.

Introduce independent edge vectors
\[
y_e=(BX)_e.
\]
Then write
\[
\min_{X,Y}
\quad
E_\ell(Y)
+
E_\Omega(Y)
+
E_R(X)
\]
subject to
\[
Y=BX.
\]
The nonlinear parts then live on \textbf{edges and vertex stars}, while the
only global consistency condition is linear.

An ADMM step gives something resembling
\[
X^{k+1}
=
\arg\min_X
E_R(X)
+
\frac{\rho}{2}
\|BX-Y^k+U^k\|^2.
\]
Ignoring or linearizing repulsion, this is
\[
\boxed{
\rho L X
=
\rho B^\top(Y-U)+F_R.
}
\]
Meanwhile the \(Y\) update decomposes into local geometric problems.

For edge lengths, it's simply sphere projection.

For solid angle, you can make a copy of each half-edge for the star around
\(v\), and the solid-angle optimization becomes a tiny independent
optimization per stitch.

That's extremely parallelizable on GPU:
\[
\text{edge projections}
\quad+\quad
\text{vertex-star projections}
\quad+\quad
\text{one global Laplacian solve}.
\]
This is very close in spirit to projective dynamics.
\href{https://users.cs.utah.edu/~ladislav/bouaziz14projective/bouaziz14projective.pdf?utm_source=chatgpt.com}{[Users at Utah]}

\bigskip\hrule

\section*{10. If edge length really is a constraint, don't optimize it as a penalty}

This may be an even bigger improvement.

Your scientific question sounds like:
\begin{quote}
Given this graph and these prescribed yarn/stitch lengths, what shape does
the graph naturally occupy?
\end{quote}
If the lengths are genuinely prescribed, then mathematically the clean
problem is
\[
\boxed{
\min_X
E_\Omega(X)+\lambda_R E_R(X)
}
\]
subject to
\[
\boxed{
\|x_i-x_j\|=\ell_{ij}
\qquad\forall(i,j)\in E.
}
\]
Then you're optimizing \textbf{on the space of realizations of the
bar-and-joint framework}.

Define
\[
c_e(X)
=
\frac12
\left(
\|x_i-x_j\|^2-\ell_e^2
\right).
\]
Its Jacobian is exactly the rigidity matrix
\[
R(X).
\]
An SQP iteration solves a KKT system of the form
\[
\boxed{
\begin{bmatrix}
H & R^\top\\
R & 0
\end{bmatrix}
\begin{bmatrix}
\Delta x\\
\lambda
\end{bmatrix}
=
-
\begin{bmatrix}
g\\
c
\end{bmatrix}.
}
\]
Now your optimizer moves principally along \textbf{flexes that preserve
yarn length}, rather than constantly stretching edges and then paying a
penalty to shrink them back.

That could be an enormous conditioning improvement if your current
edge-loss coefficient is large.

The nullspace
\[
R(X)\Delta x=0
\]
is literally the space of infinitesimal length-preserving motions of the
knit.

That is a very attractive definition of the ``available deformation
space'' of knitted fabric.

\bigskip\hrule

\section*{11. The Ricci-flow connection is real---but there is an important distinction}

Standard Ricci flow changes a \textbf{metric}, not directly an embedding.

For a triangulated piecewise-flat surface, discrete Gaussian curvature is
\[
K_v
=
2\pi-\sum_{f\ni v}\theta_{vf}.
\]
Discrete conformal Ricci flow then introduces some metric/conformal
parameter \(u_v\) and evolves it according to something like
\[
\dot u_v=K_v^\star-K_v.
\]
Importantly, the prescribed curvature \(K^\star\) does \textbf{not} need to
be zero or ``average flatness.'' Discrete Ricci-flow methods explicitly
handle user-prescribed curvature.
\href{https://sites.math.rutgers.edu/~fluo/mpapers/combinatorial%20Ricci%20flow%20in%20dimension%202.pdf?utm_source=chatgpt.com}{[Department of Mathematics]}

And there is a beautiful connection to your proposed acceleration:
\[
K(u+\Delta u)
\simeq
K(u)+J_K\Delta u.
\]
Instead of explicit Ricci flow
\[
u^{k+1}=u^k+\eta(K^\star-K),
\]
take the Newton step
\[
\boxed{
J_K\Delta u
=
K^\star-K.
}
\]
In discrete conformal geometry, \(J_K\) is Laplacian-like; the discrete
Ricci energy has a variational formulation and a Laplacian-like Hessian.
\href{https://pubmed.ncbi.nlm.nih.gov/18599915/?utm_source=chatgpt.com}{[PubMed]}

So one interpretation of what we're doing is:
\begin{quote}
\textbf{replace explicit curvature flow by an implicit/Newton curvature
solve.}
\end{quote}
That is exactly the philosophical transition that would turn your 10,000
gradual iterations into tens of large, globally coordinated changes.

\bigskip\hrule

\subsection*{But your solid angle is not automatically Ricci curvature}

This is worth being precise about.

If your solid angle \(\Omega_v\) comes from the 3-D directions of the
incident graph edges, then it is an \textbf{extrinsic quantity of the
embedding}.

Classical Gaussian/Ricci curvature is intrinsic.

So for a bare graph,
\[
\Omega_v
\]
is not literally discrete Ricci curvature.

I'd call
\[
C_v=\Omega_v-\Omega_v^\star
\]
a \textbf{vertex curvature-like residual}.

Then your flow
\[
\dot X
=
-
P^{-1}
J_\Omega^\top C
\]
is more accurately a \textbf{preconditioned curvature flow} or
\textbf{Sobolev curvature flow}.

But if your knitting graph has a natural 2-complex---e.g.\ the yarn/stitch
topology identifies actual fabric cells---things get much closer.

Triangulate those cells, calculate
\[
K_v=2\pi-\sum_f\theta_{vf},
\]
and prescribe a topology/stitch-dependent target
\[
K_v^\star.
\]
Then you really do have a discrete prescribed-curvature problem.

And rather than forcing everything toward
\[
K_v^\star=0,
\]
you could have
\[
K_v^\star
=
K(\text{knit}),
\quad
K(\text{purl}),
\quad
K(\text{increase}),
\quad
K(\text{decrease}),
\ldots
\]
or derive \(K_v^\star\) from your solid-angle model.

That would be a very interesting knitting-specific analogue of discrete
Ricci flow.

\bigskip\hrule

\section*{12. Where \(A^k\) and Laplacian eigenvectors actually help}

I would use the graph hierarchy in four specific ways rather than
explicitly putting arbitrary powers of \(A\) into the objective:
\begin{enumerate}
  \item \textbf{\(B\)} for all edge quantities.
  \item \textbf{\(L=B^\top WB\)} as your global propagation/preconditioning
        operator.
  \item \textbf{\(A^2\)} implicitly through the solid-angle Hessian, since
        vertex stars create two-hop coupling.
  \item \textbf{multiscale Laplacians} for the long-wavelength shape modes.
\end{enumerate}
That last one may be especially important.

You can build a coarse knitting graph
\[
G_0\leftarrow G_1\leftarrow\dots\leftarrow G
\]
by aggregating neighboring stitches.

Solve the shape problem on the coarse graph first, prolongate to the next
level, optimize a few times, and repeat.

That handles exactly the modes gradient descent is worst at:
long-wavelength bending.

An algebraic-multigrid hierarchy for \(L\) effectively does this
automatically.

You could also initialize from the lowest nontrivial eigenvectors
\[
Lq_i=\lambda_iq_i,
\]
but I would regard spectral embedding as an initializer, not the final
geometry. It likes smoothness and graph distance, which isn't quite the
physical objective you want.

\bigskip\hrule

\section*{13. The solver I would try first}

If I were modifying your current code, I'd implement this sequence:
\begin{enumerate}
  \item Build \(B\), \(L_0\), and an AMG or sparse-Cholesky preconditioner
        once.
  \item Replace the edge-length gradient by the local projection
        \[
        p_e=\ell_eu_e.
        \]
  \item For your log loss, calculate the lagged weights
        \[
        \alpha_e=
        w_e
        \frac{\log(r_e/\ell_e)}
        {r_e(r_e-\ell_e)}.
        \]
  \item Compute the solid-angle residual \(c\) and its sparse Jacobian
        \(J_\Omega\) with autodiff.
  \item Calculate repulsion with Barnes--Hut.
  \item Solve
        \[
        \left[
        L_\alpha\otimes I_3+
        \lambda_\Omega J_\Omega^\top J_\Omega+
        \mu L_0\otimes I_3
        \right]\Delta x=-g.
        \]
  \item Line-search the \textbf{original} energy and repeat.
\end{enumerate}
That gives you almost exactly the same objective you have now, but the
iteration has changed from a local explicit flow into a global implicit
geometric solve.

I would expect the qualitative difference to be dramatic: not ``the fabric
gradually crawls into shape,'' but rather ``each iteration makes a coherent
macroscopic deformation.'' I wouldn't promise a particular iteration count
without seeing the graph, but this is the sort of change where
\textbf{10--100 nonlinear iterations instead of \(10^4\)} is a realistic
target.

And there is a deeper conceptual formulation here that I think could be
publishable in the context of your knitting-layout work:
\[
\boxed{
\text{knitting shape}
=
\text{prescribed metric}
+
\text{prescribed vertex curvature}
+
\text{self-avoidance}.
}
\]
Edge lengths specify the discrete metric; your solid-angle term specifies a
preferred local curvature/extrinsic geometry; repulsion picks a
non-self-intersecting realization. The rigidity/Laplacian machinery then
gives the numerical method.

The particularly nice mathematical bridge to Ricci flow is that both
problems become \textbf{``compute local curvature residuals, then invert a
Laplacian-like Jacobian to make a global metric/shape correction''}, rather
than taking thousands of explicit curvature-force steps.
\href{https://pubmed.ncbi.nlm.nih.gov/18599915/?utm_source=chatgpt.com}{[PubMed]}

If you give me the exact formulas you're currently using for the
\textbf{solid-angle loss, repulsion, and edge logloss}, I can go one level
further and derive the actual \(J_\Omega\), block Hessian, and a concrete
sparse solver iteration for \emph{your exact loss}, including what can be
pre-factored and what needs to be recomputed each iteration.

\end{document}
